\begin{textsample}{NEG Dim 1 – 2011\_c5 – Score 4.00 – 2011\_c5/enade\_2011\_tec\_gest\_prod\_4e6bad5ee2.txt}  \label{ex:f1_neg_314}
CAPÍTULO 5 DISTRIBUIÇÃO DOS CONCEITOS 5.1 PANORAMA NACIONAL DA DISTRIBUIÇÃO DOS CONCEITOS A Tabela 5.1 apresenta a quantidade e distribuição de cursos de Tecnologia em Gestão da Produção Industrial participantes do ENADE/2011, por faixa de conceito e Grande Região.
A diferença entre os cursos tabulados neste capítulo e no capítulo 2 são os cursos sem conceito, em princípio, aqueles sem alunos concluintes que participassem da prova.
Observando -se os dados da Tabela 5.1, nota -se que, dos 41 cursos participantes, 16 ( 38,9 \% ) classificaram -se com conceito 3, o valor modal.
Este foi também o conceito modal da região Sudeste ( 46,5 \%, correspondendo a 13 cursos ).
O conceito 4 foi o segundo mais frequente em nível nacional ( 29,3 \%, correspondendo a 12 cursos ) e o conceito 2, o terceiro ( 17,1 \%, correspondendo a sete cursos ).
Houve, ainda, quatro cursos ( 9,8 \% ) que receberam conceito 5 e outros dois cursos ( 4,9 \% ) que receberam conceito 1.
Nenhum dos 41 cursos de Tecnologia em Gestão da Produção Industrial ficou sem conceito ( SC ).
A região Norte participou com apenas um curso, ou 2,4 \% do total nacional, que recebeu conceito 4.
A região Nordeste também participou com apenas um curso, que, como o curso da região Norte, obteve conceito 4.
Observando a tabela 5.1, nota -se que a região Sudeste apresentou o maior percentual de cursos avaliados na Área de Tecnologia em Gestão da Produção Industrial, 28 cursos, correspondendo a 68,3 \%.
Dentre eles 13 cursos obtiveram conceito 3, o conceito modal.
O conceito 1 foi atribuído a dois cursos ( 7,1 \% ) e o conceito 2 a cinco ( 17,8 \% ).
Seis cursos ( 21,4 \% ) receberam o conceito 4 e dois outros ( 7,1 \% ) receberam o conceito 5.
A região Sul contou com 11 cursos, 26,8 \% do total nacional.
A predominância do conceito 4 foi de 36,3 \%, correspondentes a quatro dos cursos da região Sul, o modal da região.
O conceito 2 foi atribuído a dois cursos ( 18,1 \% ) e o conceito 3 a três cursos ( 27,2 \% ).
O conceito 5 foi alcançado por dois cursos ( 18,1 \% ) e nenhum curso da região recebeu conceito 1.
Na região Centro-Oeste não houve participação de curso da Área de Tecnologia em Gestão da Produção Industrial. 5.2 CONCEITOS POR CATEGORIA ADMINISTRATIVA E POR GRANDE REGIÃO A Tabela 5.2 apresenta a distribuição dos cursos participantes do ENADE/2011 de Tecnologia em Gestão da Produção Industrial, por Categoria Administrativa, de acordo com os conceitos por eles alcançados, segundo as Grandes Regiões brasileiras.
Dos 41 cursos participantes, cinco ( 12,2 \% ) eram ministrados em instituições públicas e 36 ( 87,8 \% ), em privadas.
Observa -se, que as instituições privadas concentraram a maioria dos cursos.
Dos 36 cursos oferecidos pela rede privada, 15 foram avaliados com conceito 3, o modal nacional.
Os demais receberam os conceitos : 1 ( dois cursos ), 2 ( sete cursos ), 4 ( 11 cursos ) e 5 ( um curso ).
As instituições públicas participaram com cinco cursos ( 12,2 \% ).
Três deles receberam conceito 5, o conceito modal para esta Categoria Administrativa, e os outros dois receberam conceitos 3 e 4.
Na análise por região, observa -se que, o único curso participante da região Norte, que como já comentado recebeu conceito 4, era de IES privada.
A região Nordeste que também foi representada apenas por um curso da rede privada, que obteve conceito 4.
Na região Sudeste, a proporção de cursos da rede privada, 85,7 \%, foi maior do que das outras duas regiões brasileiras com curso na Área de Tecnologia em Gestão da Produção Industrial, correspondendo a 24 dos 28 cursos participantes.
Nesta categoria, na região Sudeste, o conceito modal foi 3 ( 50,0 \%, correspondendo a 12 cursos ).
Os demais foram avaliados com conceitos : 1 ( 8,3 \%, correspondendo a dois cursos ), 2 ( 20,8 \%, cinco cursos ), 4 ( 20,8 \%, cinco cursos ).
Nenhum curso da rede privada no Sudeste alcançou conceito 5.
Entre os quatro cursos oferecidos em instituições públicas na região Sudeste, a categoria modal foi o conceito 5 ( dois cursos ).
Os outros dois foram avaliados com os conceitos 3 e 4.
Na região Sul, dos 11 cursos avaliados a maioria era de IES privada.
Os conceitos recebidos por estes cursos foram : 2 ( 20,0 \%, dois cursos ), 3 ( 30,0 \%, três cursos ), 4 ( 40,0 \%, quatro cursos ) e 5 ( 10,0 \%, um curso ).
Apenas um curso de IES pública participou na região Sul e obteve o conceito 5, o máximo para este tipo de avaliação. 5.3 CONCEITOS POR ORGANIZAÇÃO ACADÊMICA E POR GRANDE REGIÃO Na Tabela 5.3 encontra -se a distribuição dos conceitos atribuídos aos cursos participantes do ENADE/2011 na Área de Tecnologia em Gestão da Produção Industrial por Organização Acadêmica, segundo as Grandes Regiões brasileiras.
Dos 41 cursos participantes, 15 eram oferecidos em Universidades, seis em Centros Universitários e os demais 20 em Faculdades.
Esta distribuição corresponde a, respectivamente, 36,6 \%, 14,6 \% e 48,8 \% dos cursos.
De acordo com os dados apresentados, dos quatro cursos avaliados com conceito 5, três eram vinculados a Universidades.
Esse tipo de Organização Acadêmica teve o conceito 3 como modal, com sete cursos ( 46,7 \% ).
Os demais cursos avaliados receberam os conceitos 2 ( 13,3 \%, dois cursos ), 4 ( 20,0 \%, três cursos ) e 5 ( também com três cursos, como já mencionado ).
Entre os cursos oferecidos em Centros Universitários, o conceito modal foi 4, com quatro cursos ( 66,7 \% ).
Os outros dois cursos, neste tipo de Organização Acadêmica, foram avaliados com conceitos 2 e 3 ( 16,7 \% \textbf{cada} um ).
Nas Faculdades, dos 20 cursos participantes, 40,0 \% ( oito cursos ) recebeu o conceito modal 3.
Dos demais cursos neste tipo de Organização Acadêmica, dois receberam o conceito 1 ( 10,0 \% ), quatro ( 20,0 \% ) obtiveram o conceito 2, cinco ( 25,0 \% ) foram avaliados com conceito 4 e apenas um ( 5,0 \% ) alcançou o conceito 5.
Considerando -se separadamente as regiões brasileiras, verifica -se que o único curso da região Norte, que recebeu conceito 4, era oferecido por uma Faculdade e, como já comentado, esta era uma instituição privada.
Este mesmo resultado foi observado na região Nordeste, ou seja, um curso oferecido por uma Faculdade privada que obteve conceito 4.
Dos 28 cursos da região Sudeste 12 eram oferecidos por Universidades.
Para estes cursos os conceitos foram : 2 ( dois cursos ), 3 ( seis cursos, conceito modal ), 4 ( dois cursos ) e 5 ( dois cursos ).
Os Centros Universitários da região Sudeste participaram com cinco cursos, dos quais três obtiveram conceito modal, 4.
Os demais receberam os conceitos : 2 ( um curso ) e 3 ( um curso ).
Os onze cursos de Faculdades foram distribuídos da seguinte forma : conceito 3 ( seis cursos, conceito modal ), conceito 1 ( dois cursos ) e 2 ( dois cursos ) \textbf{cada}, e conceito 4 ( um curso ).
Dos onze cursos da região Sul, três eram de Universidades, sendo distribuídos nos conceitos 3, 4 e 5, com um curso \textbf{cada}.
Apenas um curso era de Centro Universitário na região Sul e recebeu conceito 4.
Em Faculdades, os sete cursos participantes foram distribuídos da seguinte forma : conceito 2 ( dois cursos ), conceito 3 ( dois cursos ), conceito 4 ( dois cursos ) e conceito 5 ( um curso ).
A região Centro-Oeste, como já comentado, não foi representada com curso na Área de Tecnologia em Gestão da Produção Industrial no ENADE/2011.

% matched lemmas: cada
\end{textsample}
